% Embedding bridge between Rayleigh and RKHS operators (C1)
\subsection{Embedding bridge: Rayleigh vs.\ RKHS}\label{subsec:a3-embedding-bridge}
We fix a single heat scale $t$ and work with the Rayleigh operator
$T_P^{\mathrm{Ray}}(t,M)$ on $\CC^{2M+1}$ and the RKHS operator
$T_P^{\mathrm{RKHS}}(t)$ on $\mathcal H_t$.  The comparison
\[
  \bigl\|T_P^{\mathrm{Ray}}(t,M)\bigr\| \ \le\ \bigl\|T_P^{\mathrm{RKHS}}(t)\bigr\|
\]
is enforced by an explicit embedding $\iota_{t,M}$ and the compression formula
\[
  T_P^{\mathrm{Ray}}(t,M)\ :=\ \iota_{t,M}^\ast\,T_P^{\mathrm{RKHS}}(t)\,\iota_{t,M}.
\]
We record two concrete constructions of $\iota_{t,M}$ that each yield the same
compression inequality; the choice is driven by formalization cost.

\paragraph{C1 micro-lemmas (shared).}
\begin{enumerate}
  \item \textbf{Spaces:} set $E:=\CC^{2M+1}$ and $H:=\mathcal H_t$ as Hilbert spaces.
  \item \textbf{Embedding:} construct a linear isometry $\iota:E\hookrightarrow H$.
  \item \textbf{Compression:} define $T_P^{\mathrm{Ray}}:=\iota^\ast T_P^{\mathrm{RKHS}}\iota$.
  \item \textbf{Isometry norm:} $\|\iota\|=\|\iota^\ast\|=1$.
  \item \textbf{Submultiplicativity:} $\|AB\|\le\|A\|\,\|B\|$.
  \item \textbf{Conclusion:} $\|T_P^{\mathrm{Ray}}\|\le\|T_P^{\mathrm{RKHS}}\|$.
\end{enumerate}

\paragraph{Option A: analytic feature-map embedding.}
Choose orthonormal feature vectors $\{\psi_k\}_{k=-M}^M\subset\mathcal H_t$
compatible with the heat kernel (Fourier--Gaussian features).  Define
\[
  \iota_{t,M}(v)\ :=\ \sum_{k=-M}^{M} v_k\,\psi_k,
\]
so $\iota_{t,M}$ is a linear isometry.  The compression matrix entries satisfy
\[
  \langle T_P^{\mathrm{RKHS}}(t)\psi_j,\psi_i\rangle
  \;=\;\sum_n w_{\mathrm{RKHS}}(n)\,\Phi_{B,t}(\xi_n)\,
  \overline{\psi_j(\xi_n)}\,\psi_i(\xi_n),
\]
which reproduces the Rayleigh matrix once the feature evaluations are matched
to the trigonometric packet basis.  This is the mathematically clean bridge,
but it requires explicit RKHS feature evaluations in the formal proof.

\paragraph{Option B: pragmatic dictionary compression.}
Let $\{d_1,\dots,d_N\}\subset\mathcal H_t$ be a finite dictionary (e.g. kernel
sections $k_t(\cdot,\xi_n)$ at the prime nodes) and set
$V:=\mathrm{span}\{d_i\}$.  Choose an orthonormal basis
$\{e_1,\dots,e_N\}$ of $V$ and define the isometric inclusion
\[
  \iota_{t,M}:\CC^N\to \mathcal H_t,\qquad \iota_{t,M}(c):=\sum_{j=1}^{N} c_j\,e_j.
\]
Then $\iota_{t,M}$ is a linear isometry and the compression
$\iota_{t,M}^\ast T_P^{\mathrm{RKHS}}(t)\iota_{t,M}$ yields a finite-dimensional
operator dominated by $\|T_P^{\mathrm{RKHS}}(t)\|$.  This route avoids explicit
kernel feature computations and is typically Lean-friendly.
